\section{Conditional probabilities at intermediate resolutions}
\label{sec:conditioning}

We now pass from the measures $\mu_i$ to conditional probabilities along
unipotent plaques. The conditioning must be performed before taking the
limit: exact conditional measures of a compact diagonal orbit along a
unipotent group contain no information about its distribution at the
scales of the arithmetic estimate. We instead condition on small
\emph{transverse cells}. The freedom to choose their depth over a broad
range will make the resulting data compatible with changes of charts and
with certain diagonal motions whose length tends to infinity.

Throughout this section, $G=\SL_n(\R)$, $X=\SL_n(\Z)\backslash G$, and
$n=6$ in the application. We assume only that $\mu_i$ are $A$-invariant
probability measures converging weakly to an $A$-invariant probability
$\mu$, and that the available times satisfy
\[
 T_{-,i}\longrightarrow\infty,
 \qquad T_{+,i}/T_{-,i}\longrightarrow\infty.
\]
The particular arithmetic values of these times will be used later.
All subsequences taken below preserve these assumptions.

\subsection{Plaque charts and their extensions}

Write $L$ for the traceless diagonal algebra and $a(b)=\exp(b)$ for
$b\in L$. For $j\ne k$, set
\[
 \beta_{jk}(b)=b_k-b_j,
 \qquad U_{jk}=\{1+tE_{jk}:t\in\R\}.
\]
These signs correspond to the right action on $X$: the displacement $u$
at $x$ becomes $a(-b)ua(b)$ at $xa(b)$, and its $jk$ parameter is
multiplied by $e^{\beta_{jk}(b)}$.

A \emph{pattern} is a set of roots positive on some element of $L$ and
closed under root addition whenever the sum is a root. Its
\emph{pattern group} $V$ consists of the unipotent matrices with arbitrary
entries in those positions and zero entries elsewhere. After permuting
coordinates it is a triangular unipotent group. All these groups are
normalized by $A$, and there are only finitely many of them. We also use
the groups $VA$, but only in local charts.

For any such group $H$, a plaque chart is an injective smooth map
\begin{equation}
 \Phi:Q\times B\longrightarrow X,
 \qquad \Phi(q,h)=\psi(q)h,
 \label{eq:plaque-chart}
\end{equation}
where $Q$ is an open box in a complementary transverse manifold and
$B$ is an open bounded coordinate box in $H$. The map extends smoothly
and injectively to a slightly larger domain. The sets with $q$ fixed
are its plaques. Smooth reparametrizations of $h$, allowed to depend
smoothly on $q$, describe the same plaques and will also be used. A
smaller chart with closure in the displayed chart supplies a fixed
margin from its boundary.

The comparison of probabilities on different bounded portions of a
plaque requires charts extending beyond a prescribed displacement
range. The following elementary recurrence observation supplies them.

\begin{lemma}[No periods in a pattern group]\label{lem:no-periods}
Let $\nu$ be an $A$-invariant probability on $X$ and let $V$ be a
pattern group. For $\nu$-almost every $x$,
\[
 xv=x,\quad v\in V \quad\Longrightarrow\quad v=1.
\]
At every such $x$, the image $xK$ of any prescribed compact
displacement range $K\subset V$ is contained in an embedded plaque
with a transverse neighborhood. Consequently, for any compact displacement range in $V$
and any $\epsilon>0$, finitely many inner plaque charts cover a set of
$\nu$-measure at least $1-\epsilon$, and their outer charts contain that
range about every point of the corresponding inner charts.
\end{lemma}

\begin{proof}
Choose $b\in L$ such that $a(-tb)va(tb)\to1$ as $t\to\infty$ for every
$v\in V$. At a point recurrent under $a(b)$, a relation $xv=x$ gives
\[
 xa(kb)\,a(-kb)va(kb)=xa(kb).
\]
A subsequence of these points returns to a compact neighborhood of
$x$. Stabilizers of points in that neighborhood contain no nonidentity
element in a fixed neighborhood of $1$: this follows from the
discreteness of the lattice and local quotient charts. The displayed
relation is therefore impossible for $v\ne1$. Poincar\'e recurrence
applies at almost every point.

The orbit map $v\mapsto xv$ is then an injective immersion. Its
restriction to a compact set is an embedding. Take a slightly larger
compact coordinate domain in $V$ and use a complementary transverse
slice. Compactness permits shrinking this slice until its product
with the prescribed plaque domain is injective. Indeed, a failure
along slices shrinking to $x$ would give either two distinct points
of the limiting embedded plaque with the same image, or contradict
local injectivity near one of its points. Shrinking once more provides
inner charts with margins. A countable subcover and continuity from
below give the finite-cover assertion.
\end{proof}

Fix a countable stock of plaque charts and inner charts with the
following properties. It is a basis wherever local charts are needed;
for pattern groups it also supplies the extensions in
Lemma~\ref{lem:no-periods}, for each compact displacement range in a
countable exhaustion. It includes subboxes, fixed subdivisions of the
plaque coordinates into arbitrarily fine bins, and the simultaneous
coordinate usages described next. Chart and bin boundaries have
$\mu$-measure zero. Such choices are possible by varying coordinate
endpoints; for each coordinate only countably many levels can have
positive measure. One may start with multiplication charts and choose
countably many centers, transversals, and domain sizes dense among
the choices with compact closure and an injective enlargement.

Two simultaneous usages will matter. If $V=WU$, where $U$ is a root
group, $W$ is a normal pattern subgroup, and the patterns are disjoint,
include charts
\[
 (q,w,u)\longmapsto\psi(q)wu.
\]
Their $V$, $U$, and $W$ transverse coordinates are, respectively,
$q$, $(q,w)$, and $(q,u)$. Normality gives the last assertion; the
$w$ coordinate is a smoothly reparametrized $W$ coordinate when $u$
is fixed. Also include charts $(q,v,a)\mapsto\psi(q)va$ for $VA$, with
the simultaneous $V$ transverse coordinate $(q,a)$. All chart bounds,
plaque extensions, bin sizes, and fixed margins in a comparison are
chosen before $i$ tends to infinity.

\subsection{A broad range of almost unchanged conditionals}

In a chart $\mathcal B=\Phi(Q\times B)$, restrict $\mu_i$ to
$\mathcal B$. Partition $Q$ by a shifted dyadic grid of side $2^{-l}$.
For a point $x=\Phi(q,h)$ in the chart, let
\[
 P^{\mathcal B}_{i,l}(x)
   =\text{the conditional distribution of $h$ given the grid cell of $q$}.
\]
The conditioning uses the restricted measure. Zero-mass columns may
be assigned any probability. An expectation involving this chart
means integration over $\mathcal B$ with the unnormalized measure
$\mu_i|_{\mathcal B}$; equivalently, one can normalize that measure and
then multiply the expectation by $\mu_i(\mathcal B)$.

For positive sequences $r_i,s_i$, write $r\prec s$ if $r_i/s_i\to0$.
An interior depth will mean a sequence $d_i$ with
$l_-\prec d\prec l_+$, where the endpoints are supplied by the next
lemma. Indices $i$ on scales will often be suppressed.

\begin{lemma}[A common plateau]\label{lem:plateau}
After passage to a subsequence there are integer depths $l_{-,i}$ and
$l_{+,i}$ such that
\begin{equation}\label{eq:plateau-room}
 T_-\prec l_-\prec l_+\prec T_+.
\end{equation}
For every chart and every fixed finite plaque-bin partition in the
countable stock, the expected $\ell^1$ distance between the
conditional bin-probability vectors at depths $l_-$ and $l_+$ tends
to zero, averaged also over a uniform shift of the nested dyadic
grid. The same conclusion holds between any two intervening depths
for that grid.

The shifts can be selected so that these conclusions hold for every
fixed chart and fixed bin partition and so that random-grid boundary
estimates hold at the two endpoint depths and at any prescribed
countable collection of interior depth sequences. When coordinate
usages share transverse labels,
their grids may share the corresponding shifts. At interior depths,
conditional bin probabilities are unchanged in expected $\ell^1$
distance by using auxiliary shifted grids selected with the same
boundary convention.
\end{lemma}

\begin{proof}
Fix a chart, a finite bin partition, and a grid shift. If $p_l$ is the
conditional bin-probability vector, its coordinates form bounded
martingales in the increasing transverse grid sigma-fields. Thus
\[
 \int_{\mathcal B}\|p_v-p_u\|_2^2\,\dd\mu_i
 =\int_{\mathcal B}\|p_v\|_2^2\,\dd\mu_i
  -\int_{\mathcal B}\|p_u\|_2^2\,\dd\mu_i,
 \qquad u\le v.
\]
The sum of these differences over disjoint depth intervals is at most
$1$. This remains true after averaging in the shift. Enumerate all
fixed chart/bin choices and assign them positive summable weights.
Their weighted sum has the same bounded total variation.

Put $L_i=\log(T_{+,i}/T_{-,i})$. Divide the interval of logarithmic
depths into $m_i$ consecutive intervals, where $m_i\to\infty$ and
$m_i=o(L_i)$. Discard its first and last intervals. Some remaining
interval has weighted martingale cost $O(m_i^{-1})$. Its endpoints,
rounded to integers, give $l_-$ and $l_+$. The discarded end intervals
and the length $L_i/m_i\to\infty$ give all three ratios in
\eqref{eq:plateau-room}. Each fixed chart/bin choice has cost tending
to zero; Cauchy--Schwarz gives the asserted $\ell^1$ conclusion.
Orthogonality of martingale increments gives it for intervening
depths as well.

For the boundary assertion, fix a point $q$ in a Euclidean transverse
coordinate domain. Under a uniform grid shift, the probability that
$q$ is within $\rho 2^{-l}$ of a depth-$l$ grid wall is $O(\rho)$,
uniformly in $q,l$, and the sampling measure. Integrating first in
the shift gives the same estimate for any measure on that domain.
Choose $\rho_i\downarrow0$ slowly. Markov's inequality, followed by
a diagonal choice over the chart/bin and depth lists, selects shifts
with both the conditional-probability estimates and boundary masses
tending to zero. Common labels can use common shift variables:
each marginal shift remains uniform, which is all the argument uses.

Here is the comparison with an auxiliary grid at an interior depth
$d$. Outside negligible boundary mass, one of its cells lies inside
one depth-$l_-$ cell and is a union of depth-$l_+$ cells, apart from
the latter cells meeting its walls. The relevant relative collar
widths are $O(2^{l_--d})$ and $O(2^{d-l_+})$. Both tend to zero much
faster than the chosen $\rho_i$. Summing the martingale bin errors
over the contained fine cells compares the auxiliary conditional
law to the coarse one. The discarded mass changes the weighted
conditional $\ell^1$ error by at most twice that mass. This proves
the last assertion. The same argument compares any two of the
selected grids by sandwiching them between these endpoint grids.
\end{proof}

We make the selection convention precise for later use. All defining
charts, simultaneous coordinate usages, fixed bin subdivisions, and
depth sequences belong to countable lists. A diagonal selection
enforces their estimates for each fixed member; it does not assert
uniformity over the whole infinite list at a finite index $i$.
Within $[l_-,l_+]$ any fixed finite number of scales can be chosen
with successive ratios tending to infinity and with room at both
ends. For instance, use distinct fixed-fraction interpolations on
the logarithmic depth interval. Different arguments may reuse such
scales. Auxiliary grids chosen for one comparison need not define
the data, because Lemma~\ref{lem:plateau} compares their fixed-bin
laws with those of the defining grids. Such a grid may also be
chosen later at an interior depth not in the original list; it
does not change the augmented states defined below. We always let $\rho_i$ decrease
slowly enough to dominate the exponentially small relative collar
widths arising from these scale separations. None of these conventions
requires absolute continuity.

\subsection{The augmented probability space}

Choose a defining interior depth for every chart and selected grids
as above. Record at $x$ its conditional plaque probability in each
chart containing it, and an isolated dummy symbol for every other
chart. Write $Z_i(x)$ for the resulting tuple together with $x$.
Each plaque probability is regarded as a probability on the compact
closure of its coordinate domain. The countable product of these
compact spaces, together with the $X$ coordinate, is a separable
metrizable space. Tightness of $\mu_i$ therefore permits passage to
a subsequence for which the laws of $Z_i(x)$ converge. Denote the
limiting state by
\[
 \omega=\bigl(x,(P^{\mathcal B})_{\mathcal B}\bigr).
\]
Its $x$ marginal is $\mu$. Fixed-bin probabilities pass to the limit:
the averaged plaque probability in a chart is its plaque-coordinate
marginal, so the null-boundary choices imply that the limiting
probabilities assign zero mass to the chosen bin boundaries almost
surely.

\begin{lemma}[Sampling one plaque probability]\label{lem:sampling}
In any fixed chart, conditional on its transverse coordinate $q$
and its own limiting probability $P^{\mathcal B}$, the sampled
plaque coordinate $h$ has law $P^{\mathcal B}$. In particular,
$h\in\supp P^{\mathcal B}$ almost surely on that chart. This assertion
conditions on this one probability, rather than on the entire
augmented state.
\end{lemma}

\begin{proof}
Let $q_i^0$ be a representative of the cell containing $q$. At the
sequence level, $P^{\mathcal B}_{i}$ is a function of that cell.
Consequently, for bounded continuous tests $f$ and $F$,
\[
 \int_{\mathcal B} f(h)F(q_i^0,P^{\mathcal B}_{i})\,\dd\mu_i
 =\int_{\mathcal B}P^{\mathcal B}_{i}(f)
                  F(q_i^0,P^{\mathcal B}_{i})\,\dd\mu_i.
\]
On compact inner boxes, replacing $q_i^0$ by $q$ incurs an error
tending to zero. Let $i\to\infty$ and then exhaust the chart by
inner boxes. Its boundary has zero limiting mass. This proves the
conditional-law identity. The standard support assertion for a
sample from a probability measure follows, for example, by using a
countable basis of the plaque domain.
\end{proof}

The remaining task in this section is to show that probabilities
recorded in different charts describe the same measure along a
plaque, up to normalization. We first isolate the estimate that
allows many small transverse columns to be regrouped without losing
control of the error.

\subsection{Aggregating fine columns into patches}

In an output chart use the defining nested grid at depths $l_-$ and
$l_+$. Write $E$ for a coarse cell and $e$ for a fine cell contained
in $E$. Fix finitely many plaque bins $J_j$, and set
\[
 M_e=\mu_i(\text{$e$ column}),\qquad
 M_{e,j}=\mu_i(\text{$e$ column with }h\in J_j),\qquad
 P_E(j)=\Prb_i(h\in J_j\mid q\in E).
\]
Here and below the chart restriction is understood. The plateau
estimate is precisely
\begin{equation}\label{eq:mass-error}
 \delta_i:=\sum_E\sum_{e\subset E}\sum_j
          |M_{e,j}-M_eP_E(j)|\longrightarrow0.
\end{equation}

\begin{lemma}[Comparison on trimmed patches]\label{lem:patch-comparison}
In the preceding chart, let $p$ range over pairwise disjoint
measurable patches, each lying transversely in one coarse cell
$E(p)$. Suppose a subset of each patch has the product form
\[
 p^{\mathrm{in}}=
 \bigcup_{e\in I_p}\bigl(e\times B_p\bigr),
 \qquad B_p=\bigcup_{j\in J(p)}J_j,
\]
and the sum of the discarded masses $\mu_i(p\setminus p^{\mathrm{in}})$
is at most $\epsilon$. The collections $I_p,J(p)$ may depend on
$i,p$, and the number of patches need not be bounded. If $L_p$ is
the plaque-bin law conditional on $p$, then, with errors weighted by
$\mu_i(p)$,
\[
 L_p\quad\hbox{agrees with}\quad
 \frac{\one_{B_p}P_{E(p)}}{P_{E(p)}(B_p)}
 \quad\text{to total expected }\ell^1\text{ error }
 O(\epsilon)+o(1).
\]
An arbitrary law is used if the denominator vanishes. The
unnormalized version is
\begin{equation}\label{eq:patch-unnormalized}
 \sum_p\mu_i(p)
 \left\|P_{E(p)}(B_p)L_p-
             \one_{B_p}P_{E(p)}\right\|_1
 \le O(\epsilon)+o(1).
\end{equation}
If patches and a disjoint refinement into subpatches admit such trims with the same
coarse cell and the same plaque-bin union, their conditional bin
laws agree with one another with the corresponding weighted error.

These assertions also hold for bounded continuous tests, up to their
uniform continuity error on the fixed bins. A test may depend on
the sampled point, provided its norm and continuity modulus in the
tested plaque variable are uniform. At the output sample,
$P_{E(p)}$ may be replaced in fixed-bin expressions by the recorded
data probability, with an additional error tending to zero.
\end{lemma}

\begin{proof}
For a trimmed patch, sum the errors in \eqref{eq:mass-error} over
$e\in I_p$ and $j\in J(p)$. The resulting comparison is between
its unnormalized bin-mass vector and
\[
 \left(\sum_{e\in I_p}M_e\right)
             \one_{B_p}P_{E(p)}.
\]
Every pair $(e,j)$ occurs in at most one trimmed patch, since those
sets are disjoint. Hence the total unnormalized error is at most
$\delta_i$. Restoring discarded mass adds at most $\epsilon$.
For two nonnegative finite vectors $v,w$, comparison after
normalization, weighted by $\|v\|_1$, costs at most
$2\|v-w\|_1$; this follows by adding and subtracting
$w/\|v\|_1$. The same inequality covers zero mass by an arbitrary
choice of normalized vector. It proves the first assertion and,
on multiplying by $P_{E(p)}(B_p)\le1$, proves
\eqref{eq:patch-unnormalized}.

Apply the same argument separately to patches and subpatches and
use their common restricted vector $P_E|_{B_p}$ to obtain their
comparison. Approximating a continuous test by its values on bins
gives the remaining test assertion. The bound is in $\ell^1$, so it
is uniform over all choices of bounded bin values, including values
depending on the sample. Finally, the plateau controls the expected
fixed-bin distance between the coarse conditional vector and the
recorded vector. Restricting to disjoint patches does not increase
that expectation.
\end{proof}

In applications the bins and charts are fixed first. Their sizes
may then be chosen sufficiently small to achieve any prescribed
limsup error, after which $i$ tends to infinity. Thus ``arbitrarily
small error'' below means an actual bound by a prescribed
$\epsilon>0$ in this order of choices. It is not a claim about
uniform convergence over all subdivisions.

\subsection{Projective agreement and diagonal motions}

Center the probability in a $V$ chart at its sampled point: if the
sample has group coordinate $h_s$, the group displacement of $h$ is
$h_s^{-1}h$. For a smoothly reparametrized chart, use the corresponding
group elements in this formula. Two nonzero measures on a common
domain agree \emph{projectively} if one is a positive scalar multiple
of the other there.

\begin{proposition}[Comparison of plaque data]\label{prop:plaque-comparison}
For the augmented limit constructed above the following assertions
hold.
\begin{enumerate}
\item For every pattern group $V$, its centered probabilities in
different charts agree projectively on their common plaque sheet
near the sampled point. They extend to a measurable nonzero
projective locally finite measure $m^V_\omega$ on $V$, with
$1\in\supp m^V_\omega$. The chart probabilities represent its
restrictions on all bounded displacement ranges contained in their
plaque domains.
\item Fix $b\in L$. In every subsequential joint limit of
\[
 \bigl(Z_i(x),Z_i(xa(b))\bigr),\qquad x\sim\mu_i,
\]
the two $V$ measures are related projectively by the pushforward
under $v\mapsto a(-b)va(b)$.
\item Fix a root $\beta$ and an interior time scale $R$ with room
for interior depths on both sides. For every fixed $C<\infty$ and
every sequence
\[
 b_i\in\ker\beta,\qquad \|b_i\|\le CR_i,
\]
the measures $m^{U_\beta}$ in every subsequential joint limit of
$\bigl(Z_i(x),Z_i(xa(b_i))\bigr)$ agree projectively. This assertion
holds for every such sequence $b_i$; the defining grids and state
construction are unchanged.
\end{enumerate}
\end{proposition}

\begin{proof}
We prove one local comparison that covers all three assertions. The
motion is either the identity, a fixed $a(b)$, or $a(b_i)$ with
$b_i\in\ker\beta$ and $V=U_\beta$. Denote its induced conjugation on
$V$ by $c$. In the first two cases $c$ is fixed; in the third case
$c$ is the identity.

We display the geometry in group-element plaque coordinates. For
a reparametrized chart, write its original coordinate as $t\in B_0$
and its group element as $h=r_q(t)$, where the maps $r_q$ and their
inverses are smooth on the fixed extended chart. Its group-coordinate
domain is $B(q)=r_q(B_0)$; for a multiplication chart $B(q)$ is
constant. The conditional probabilities and bins remain in the
original fixed coordinates. When converting a probability to group
coordinates at a sample, use $r_{q_s}$. Across a depth-$F$ cell this
conversion differs from the actual $r_q$ by $O(2^{-F})$, which has
vanishing effect on every continuous test considered below.

In the long-motion case choose an input
transverse depth $F$ with
\begin{equation}\label{eq:comparison-scales}
 l_-\prec R\prec F\prec l_+,
\end{equation}
For a fixed motion, it suffices to take $l_-\prec F\prec l_+$;
all factors $\exp(O(R))$ below are then bounded constants.
The depth-$F$ input grid may be an auxiliary grid selected for this
comparison. Thus $R$ and $F$ need not belong to the countable list
used to define the states; Lemma~\ref{lem:plateau} compares the
auxiliary laws back to the recorded ones. Its boundary selection
uses $\mu_i$ and the scale bound, and does not depend on the
particular $b_i$ within that bound.
The input patches are the depth-$F$ transverse cells times the full
plaque-coordinate domain. Their conditional laws are the input
probabilities at depth $F$.

\smallskip\noindent
\emph{Fitting the images into output charts.}
Choose input inner charts and fixed displacement test ranges first.
An input chart may have any fixed plaque extension needed for these
ranges. By Lemma~\ref{lem:no-periods}, output inner charts can be
chosen whose outer charts accommodate the entire image plaque
range, with a margin. A finite such cover omits arbitrarily small
limiting mass. The output distribution is again $\mu_i$, so the
same is true for the sequence. If an image patch meets one of these
inner charts, its whole image fits in the corresponding outer chart
for all sufficiently large $i$. This last assertion follows from
the transverse diameter estimate below. Assign such patches to
output charts, or work with each applicable pair of charts; the
finite number of pairs only multiplies error bounds by a fixed
constant.

More precisely, let $B(q)$ and $B'(q')$ be the full input and output
domains in group coordinates, and let $(q_s,h_s),(q'_s,h'_s)$ be the
two samples. The output extension is required to contain
\begin{equation}\label{eq:full-plaque-margin}
 h'_s\,c\bigl(h_s^{-1}\overline{B(q_s)}\bigr)\subset B'(q'_s)
 \quad\text{with a positive boundary margin}.
\end{equation}
This condition concerns the entire input plaque domain. It is an
open condition on the paired sample coordinates: the domains are
smooth images of fixed compact closures. Their possible centered
ranges lie in a fixed compact subset of $V$, so the extension lemma
supplies output covers with this property.

In the two charts the map has the form
\begin{equation}\label{eq:plaque-transition}
 (q,h)\longmapsto\bigl(\theta_i(q),\,d_i(q)c(h)\bigr).
\end{equation}
The first coordinate is independent of $h$, because the motion
normalizes $V$. The factor $d_i(q)$ is a left translation in $V$.
It remains in a fixed compact set on the comparisons in question:
both plaque-coordinate ranges and $c$ are fixed. For the varying
motion, right multiplication and its inverse have differential norm
at most $\exp(O(R))$ in a left invariant metric on $G$ and the
induced quotient metric. Fixed chart changes have bounded
differentials. Consequently both $\theta_i(q)$ and $d_i(q)$ vary over an
input cell by
\begin{equation}\label{eq:transverse-image-width}
O\bigl(2^{-F}\exp(O(R))\bigr)=o(2^{-l_-}).
\end{equation}
For reparametrized charts, the input domains $B(q)$ also vary by
$O(2^{-F})$ over the cell, and their admitted output ranges vary by
$O(2^{-F}\exp(O(R)))$. The fixed boundary margins absorb these
variations. Thus the same estimate applies to those charts.
The same bounds hold for pullbacks of paths in the extended charts.
They justify following one plaque sheet throughout its prescribed
range, rather than including a different sheet at a return of the
orbit. They also justify the fitting assertion: vary $q$ along the
short path in its input cell and then follow the fixed plaque range
inside the output extension.

On one such patch, $\theta_i$ is injective. To see this, the admitted
output plaque ranges have a common open subset, by the margin and
the small variation of $d_i(q)$. If two transverse inputs had the
same image, one could choose the same output plaque point in this
common subset, contradicting injectivity of the motion and of the
input chart. This observation permits product trims in output
coordinates, even if many distinct input patches return to the
same output box.

\smallskip\noindent
\emph{Trimming to fine columns and fixed plaque bins.}
Discard image patches meeting a coarse transverse wall. By
\eqref{eq:transverse-image-width}, every point of such a patch is in
a collar whose relative width tends to zero. The grid convention,
applied to the output distribution $\mu_i$, makes the discarded
mass $o(1)$. Each remaining patch lies transversely in one coarse
cell $E$.

Take fine transverse cells whose closures lie in its projected
transverse range, and plaque bins admitted over the whole range.
Their products form an inner set as in
Lemma~\ref{lem:patch-comparison}. Here are the two loss estimates.
First, a path varying the transverse output coordinate through one
fine cell pulls back a distance
\begin{equation}\label{eq:fine-pullback-width}
 O\bigl(2^{-l_+}\exp(O(R))\bigr)=o(2^{-F}).
\end{equation}
Away from the input cell walls and a fixed small collar of the
input chart boundary, it therefore stays in the input transverse
range and in the chart. Input grid-wall loss is $o(1)$ by the
selected shifts; chart-boundary loss can be made arbitrarily small
by its null limiting mass. This verifies that almost all sampled
fine transverse cells are admitted.

Second, on each plaque the transformation in
\eqref{eq:plaque-transition} has bounded distortion, uniformly in
$i$, and its dependence on $q$ tends to zero over the patch. A
sample away from a small collar of the input plaque boundary thus
lies in an output plaque bin admitted over the whole transverse
range, provided the fixed bins are chosen sufficiently fine. The
collar has arbitrarily small sampling mass, again by the fixed
null-boundary choice. For a reparametrized chart these collars are
pulled back to its original fixed coordinate boundaries, which
have zero limiting mass. Continuous tests expressed in the fixed
coordinates pull back with uniform continuity moduli on the compact
chart ranges.
This gives arbitrarily small total trim loss.
No bound on the number of patches is involved.

\smallskip\noindent
\emph{Comparing the conditional laws.}
The image of $\mu_i$ restricted to an input patch is exactly
$\mu_i$ restricted to its image, by $A$ invariance (or trivially for
the identity motion). Apply Lemma~\ref{lem:patch-comparison} to its
trim. The input patch law, converted by
\eqref{eq:plaque-transition}, agrees with $P_E$ restricted to the
admitted bins, up to arbitrarily small weighted error.

To express this without dividing by a possibly small admitted
mass, take two continuous displacement tests $f,\phi$ supported
with margin inside the input plaque range. Their output bin
approximations are wholly inside the admitted bins. The
unnormalized comparison says that each input test integral,
multiplied by $P_E(B_p)$, agrees with the corresponding output
$P_E$ integral. Cross-multiplication therefore gives
\begin{equation}\label{eq:projective-cross-products}
 P_{\mathrm{in}}(f)P_{\mathrm{out}}(\phi\circ c^{-1})
 -P_{\mathrm{in}}(\phi)P_{\mathrm{out}}(f\circ c^{-1})
 \longrightarrow0
\end{equation}
in probability on the compared samples. All integrals here are in
displacement coordinates centered at the respective samples.
The small $q$ variation changes the tests by a uniformly vanishing
amount; sample-dependent centering is permitted by the uniform
test assertion in Lemma~\ref{lem:patch-comparison}. The input
depth-$F$ law can be replaced by its recorded probability by
Lemma~\ref{lem:plateau}; the output coarse law can be replaced by
the output recorded probability by the same lemma. This proves
\eqref{eq:projective-cross-products} for the data that actually
define the states.

For completeness, the errors just used are localized as follows.
Fix a pair of charts and open conditions placing the samples in
inner boxes, requiring the full-plaque margin
\eqref{eq:full-plaque-margin}, and placing the displacement test
supports inside the input plaque range with positive margins.
The patch fits and trims hold for all sequence samples satisfying
these conditions, uniformly for
the chosen motion bound. Trim accuracy may be made arbitrarily
good independently of the fixed cover or these margin conditions.
Thus the cross-product discrepancy tends to zero in probability
on each such event. Exhaust by stricter margins and use continuous
tests on compact inner boxes. Every joint distributional limit
satisfies the corresponding exact identity there, since its tests
and coordinate changes are continuous. Countably many charts,
margin conditions, and tests suffice.

\smallskip\noindent
\emph{Gluing along one plaque.}
Apply this result first to the identity motion. Compare any chart
at the sample with a chart extending far enough to contain its
bounded plaque domain. By Lemma~\ref{lem:sampling}, each probability
has positive mass on every neighborhood of the sampled point.
Choose a nonnegative continuous test supported in a common such
neighborhood and positive at that point. Its integral is positive
for both probabilities. The cross-product identities then give
projective agreement for every continuous test with compact support
in the first plaque domain.

To compare two bounded domains, use a single extension containing
both, on the same plaque sheet. Exhaust $V$ by compact connected
coordinate ranges containing $1$ in their interiors. Successive
extensions agree after normalization near $1$, and so define a
locally finite measure on the union. Local finiteness follows
because each compact set lies in one extension carrying a finite
chart probability with a finite positive normalization factor.
The identity lies in its support by the sampling lemma. This
construction does not require chaining through intermediate support
points: each comparison uses a chart containing both the identity
neighborhood and the entire displacement range in question.
Using countably many tests and choosing normalizing bumps from a
countable collection makes the resulting projective measure
measurable.

\smallskip\noindent
\emph{Equivariance, including the long motion.}
For a fixed diagonal motion, the same cross-product identities
give pushforward by $c(v)=a(-b)va(b)$, first on bounded test ranges
and then globally by the extensions just constructed. For the
varying motion in $\ker\beta$, $c$ is the identity, and all estimates
above are uniform for $\|b_i\|\le CR_i$: the only growing chart cost
is $\exp(O(R_i))$, dominated by
\eqref{eq:comparison-scales}. In a paired-state limit the second
base point need not be a fixed function of the first. This causes
no difficulty. The local identities hold on every paired open
margin event, and these events exhaust probability one because
both marginals admit the required extended charts. Their exact
identities therefore hold in every such limit. Since the argument
applies to an arbitrary sequence $b_i$ under the displayed bound,
it proves the last assertion with the stated quantifiers.
\end{proof}

The measures $m^V_\omega$ retain information obtained by conditioning
the original sequence, while having the chart consistency and
diagonal comparison properties of ordinary plaque measures. The
next step is to relate the data for different pattern groups;
that requires a product argument, rather than a further change of
coordinates.
